% GENERATED FILE. Edit canonical lessons, then run npm run generate:books.
% canonical-source-hash: fnv1a64:44b53c02343eb2a0
% canonical-lessons: FR-C01-salut, FR-W01-salut-observe, FR-W01-salut-guided-copy, FR-W01-salut-delayed-copy, FR-W01-salut-dictation, FR-C01-bien, FR-C01-bon, FR-C01-le-la, FR-C01-jour, FR-C01-bonjour, FR-C01-soir, FR-C01-bonsoir, FR-C01-nuit, FR-C01-bonne-nuit, FR-C01-practice

\chapter{Greetings}
\label{ch:greetings}
\hlchaptermodality{1}

\begin{chapteropening}
\textbf{By the end of this chapter:} \emph{I can greet someone in French, recognize the reply, and choose an appropriate response.}

The French greetings, built from their pieces — a word for “good,” words for “day,” “evening” and “night,” and the gender machine that decides their endings. French and Spanish are both daughters of Latin; where a French word and its Spanish cousin split from one Latin root, both appear here, because the difference between them is often the lesson — and the text supplies the Spanish word whether or not you have met it before.
\end{chapteropening}

\section[salut]{salut — \textquotedblleft{}hi,\textquotedblright{} and a wish for your health}

\label{lesson:FR-C01-salut}



\begin{quote}
Your first French word, and a warmer one than it looks. \emph{Salut} is the casual \textquotedblleft{}hi\textquotedblright{} — but taken apart, it literally wishes you good health.
\end{quote}

\begin{sounds}
\begin{itemize}
\raggedright
  \item \texttt{silent-final} — the final \textbf{t is silent}: \emph{salut} is \emph{sa-LU}, not \emph{sa-LUT}. Most French final consonants go silent. $\to$ reference
  \item \texttt{vowel-u} — the \textbf{u} is the French front-rounded \emph{ee} (say \emph{ee} with pursed lips), not \emph{oo}.
\end{itemize}
\end{sounds}

\begin{cousinweb}
\textbf{salut} — informal \textquotedblleft{}hi\textquotedblright{} (and \textquotedblleft{}bye\textquotedblright{}). Root: Latin \textbf{salus, salutem} (\textquotedblleft{}health, safety, well-being\textquotedblright{}). So \emph{salut} is, at root, a wish of \emph{health} — the same instinct as English \textquotedblleft{}cheers\textquotedblright{} or old \textquotedblleft{}hail\textquotedblright{} (be healthy). The family you already own in English:

\begin{itemize}
\raggedright
  \item \textbf{salut}e, \textbf{salut}ation — a greeting (a wishing-of-health).
  \item \textbf{salut}ary, \textbf{salu}brious — health-giving.
  \item and, from the close cousin \emph{salvus} (\textquotedblleft{}safe\textquotedblright{}), \emph{save}, \emph{salvation}, \emph{safe}.
\end{itemize}

\begin{quote}
Spanish note: its casual \textquotedblleft{}hi\textquotedblright{} is \emph{hola} — an uncertain, rootless interjection. French \emph{salut} is the opposite: a transparent Latin word meaning \textquotedblleft{}health.\textquotedblright{} Two Romance languages, two totally different casual \textquotedblleft{}hi\textquotedblright{}s.
\end{quote}
\end{cousinweb}

\begin{culture}
\emph{Salut} is strictly \textbf{informal} — friends, family, peers — and it doubles as \textquotedblleft{}bye.\textquotedblright{} For anyone you'd show respect to, or a first meeting, you reach past it for \emph{bonjour} (coming in a few lessons), exactly as Spanish reaches past \emph{hola} for \emph{buenos días}.
\end{culture}

\subsection*{Your turn}
Say these aloud:

\begin{itemize}
\raggedright
  \item \textquotedblleft{}salut\textquotedblright{} — \emph{sa-LU}, silent \emph{t}
  \item \textquotedblleft{}salut\textquotedblright{} then English \textquotedblleft{}salute\textquotedblright{}, \textquotedblleft{}salutary\textquotedblright{} — the health root
\end{itemize}

\begin{tcolorbox}[breakable,colback=teal!4,colframe=teal!35!black,title={Before you move on}]
What does \emph{salut} literally wish someone? (Health — from Latin \emph{salus}.) What happens to its final \emph{t}? (Silent.)
\end{tcolorbox}
\section[salut]{\emph{salut} — let your eye and hand meet one known word}

\label{lesson:FR-W01-salut-observe}



\begin{quote}
Say the casual greeting you just learned: \emph{salut}, \emph{sa-LU}. Today adds no new French. You will only notice how that familiar word sits on a page.
\end{quote}

\begin{grammarlens}[title={What you've built — one spelling}]
\begin{quote}
\textbf{s a l u t} $\to$ \textbf{salut}
\end{quote}

The word has five written letters. Your voice finishes on \textbf{u}; the final \textbf{t} stays visible and silent. Point to each letter from left to right once.
\end{grammarlens}

\subsection*{Writing — trace with the model visible}
Keep \textbf{salut} visible. Slowly finger-trace the five letters once, left to right. Then trace the word once in the air. Say \emph{sa-LU} only after your finger reaches the silent \textbf{t}.

This is noticing, not a handwriting test. Use the letter forms you normally write; there is no new alphabet and no demand for perfect penmanship.

\begin{tcolorbox}[breakable,colback=teal!4,colframe=teal!35!black,title={Before you move on}]
Point to the silent letter. Stop here; copying comes in the next tiny lesson.
\end{tcolorbox}
\section[salut]{\emph{salut} — one supported copy}

\label{lesson:FR-W01-salut-guided-copy}



\begin{quote}
Point to the last letter in \textbf{salut}. It is written \textbf{t} and said as silence. Keep the model on the page for the whole lesson.
\end{quote}

\subsection*{Writing — guided copy}
Copy \textbf{salut} once. You may look back after every letter.

When you finish, touch \textbf{s-a-l-u-t} in your copy and say \emph{sa-LU}. The eye checks five letters; the voice does not add a final \emph{t} sound.

\begin{tcolorbox}[breakable,colback=teal!4,colframe=teal!35!black,title={Before you move on}]
Compare only the order and presence of the five letters. Repair one letter if needed, then stop. The model stays visible today.
\end{tcolorbox}
\section[salut]{\emph{salut} — look, hide, write, compare}

\label{lesson:FR-W01-salut-delayed-copy}



\begin{quote}
Look at \textbf{salut} for five seconds. Touch the final silent \textbf{t}. Do not copy it yet.
\end{quote}

\subsection*{Writing — delayed copy}
\begin{enumerate}
\raggedright
  \item cover the model;
  \item wait five seconds;
  \item write the word once from memory;
  \item uncover \textbf{salut}, compare, and repair only the place that differs.
\end{enumerate}

If one letter differs, circle that place and repair only that letter. A repair is part of the exercise, not a failure.

\begin{tcolorbox}[breakable,colback=teal!4,colframe=teal!35!black,title={Before you move on}]
Read your repaired word as \emph{sa-LU}. Stop after one memory attempt.
\end{tcolorbox}
\section[salut]{One heard greeting — write what you know}

\label{lesson:FR-W01-salut-dictation}



\begin{quote}
Cover the answer line lower on this page. You will hear one familiar casual greeting. No new vocabulary is hiding in this lesson.
\end{quote}

\subsection*{Writing — short dictation}
{[}YOU HEAR: \emph{sa-LU}{]}

Write the French greeting from the sound alone. Remember that your pen records one final letter that your voice does not pronounce.

Now uncover and compare: \textbf{salut}. Check only the five letters and the silent final \textbf{t}. Repair one place if needed.

\begin{tcolorbox}[breakable,colback=teal!4,colframe=teal!35!black,title={Before you move on}]
Writing from sound is the endpoint of this four-step runway. It is not yet a DILF mock; later lessons will grow from one known word to contact details, numbers, forms, and short practical messages.
\end{tcolorbox}
\section[bien]{bien — \textquotedblleft{}well,\textquotedblright{} a whole answer}

\label{lesson:FR-C01-bien}



\begin{quote}
A tiny, essential word: \textbf{bien}, \textquotedblleft{}well.\textquotedblright{} It carries a good chunk of everyday conversation, and its root reaches straight into English.
\end{quote}

\begin{sounds}
\begin{itemize}
\raggedright
  \item \texttt{nasal-in} — \textbf{ien} is a \textbf{nasal} \emph{byan}: air through the nose, and the \emph{n} itself is silent. Not \textquotedblleft{}bee-en.\textquotedblright{} $\to$ reference
\end{itemize}
\end{sounds}

\begin{cousinweb}
\textbf{bien} — \textquotedblleft{}well.\textquotedblright{} Root: Latin \textbf{bene} (\textquotedblleft{}well\textquotedblright{}). The English \emph{bene-} family is yours already:

\begin{itemize}
\raggedright
  \item \textbf{bene}fit (\textquotedblleft{}do well\textquotedblright{}), \textbf{bene}volent (\textquotedblleft{}wish well\textquotedblright{}), \textbf{bene}diction (\textquotedblleft{}speak well\textquotedblright{}), be\textbf{nign}.
\end{itemize}

\begin{quote}
A cross-language aside: Spanish, another daughter of Latin, has the very same word — \emph{bien}, also from \emph{bene}, also \textquotedblleft{}well.\textquotedblright{} French and Spanish each inherited the Latin adverb independently and it barely changed in either; when two Romance languages agree this exactly, you're looking straight at the shared Latin parent.
\end{quote}

\textbf{Using it alone}: \textquotedblleft{}\emph{Comment ça va?}\textquotedblright{} (how's it going?) $\to$ \textquotedblleft{}\textbf{Bien.}\textquotedblright{} (\textquotedblleft{}Well / good.\textquotedblright{}) Intensified: \textquotedblleft{}\textbf{très bien}\textquotedblright{} (\textquotedblleft{}very well\textquotedblright{}).
\end{cousinweb}

\begin{grammarlens}[title={Grammar Lens: the \textquotedblleft{}good\textquotedblright{} it pairs with}]
\emph{bien} is the \textbf{adverb} (\textquotedblleft{}well\textquotedblright{} — \emph{how} something is done). The \textbf{adjective} \textquotedblleft{}good\textquotedblright{} (describing a thing) is a different word, \textbf{bon} — the next lesson. This is the same \emph{well} vs. \emph{good} split as English, and the same \emph{bien} vs. \emph{bueno} split as Spanish; French's adjective is \emph{bon}, from Latin \emph{bonus}.
\end{grammarlens}

\subsection*{Your turn}
Say these aloud:

\begin{itemize}
\raggedright
  \item \textquotedblleft{}bien\textquotedblright{} — nasal \emph{byan}
  \item \textquotedblleft{}très bien\textquotedblright{}
  \item \textquotedblleft{}bien\textquotedblright{} then English \textquotedblleft{}benefit\textquotedblright{}, \textquotedblleft{}benevolent\textquotedblright{}
\end{itemize}

\begin{tcolorbox}[breakable,colback=teal!4,colframe=teal!35!black,title={Before you move on}]
French \emph{bien} and Spanish \emph{bien} come from the same Latin word — which one? (\emph{bene}, \textquotedblleft{}well.\textquotedblright{}) Adverb or adjective? (Adverb — \textquotedblleft{}good\textquotedblright{} the adjective is \emph{bon}, next.)
\end{tcolorbox}
\section[bon / bonne]{bon / bonne — \textquotedblleft{}good,\textquotedblright{} the block you build greetings from}

\label{lesson:FR-C01-bon}



\begin{quote}
This is the word that starts \emph{bonjour} and \emph{bonsoir}. Learn it alone first — including the trick that its two forms \emph{sound} completely different even though they look nearly the same.
\end{quote}

\begin{sounds}
\begin{itemize}
\raggedright
  \item \texttt{nasal-on} — \textbf{bon} is a \textbf{nasal} \emph{bõ}: the \emph{n} is silent, the vowel goes through the nose. $\to$ reference
  \item \texttt{silent-final} — but \textbf{bonne} (feminine) is \emph{bun} — the doubled \emph{nn} \textbf{cancels the nasal} and you actually pronounce an \emph{n}. So \emph{bon} (\emph{bõ}) and \emph{bonne} (\emph{bun}) sound different despite looking alike.
\end{itemize}
\end{sounds}

\begin{cousinweb}
\textbf{bon} — \textquotedblleft{}good\textquotedblright{} (adjective, masculine). \textbf{bonne} — \textquotedblleft{}good\textquotedblright{} (feminine). Root: Latin \textbf{bonus / bona} (\textquotedblleft{}good\textquotedblright{}). Your English cousins:

\begin{itemize}
\raggedright
  \item \textbf{bon}us (borrowed straight from Latin \emph{bonus}), \textbf{bon}anza, and — via Old French \emph{bonté} — \textbf{boun}ty. Even \emph{bonbon} (\textquotedblleft{}good-good,\textquotedblright{} a candy).
\end{itemize}

\begin{quote}
Spanish cousin: from the same \emph{bonus / bona}, Spanish made \emph{bueno / buena}. Notice French kept it shorter (\emph{bon}) while Spanish broke the \emph{o} into a diphthong (\emph{bue-}). Same Latin word, two accents of wear.
\end{quote}
\end{cousinweb}

\begin{grammarlens}[title={Grammar Lens: adjectives agree with their noun}]
Like every Romance adjective, \emph{bon} \textbf{agrees} with its noun's gender and number — \emph{bon} (masc.), \emph{bonne} (fem.), \emph{bons} (masc. pl.), \emph{bonnes} (fem. pl.). This is Latin's legacy: adjectives change their ending to match the noun (Spanish runs the same machine on \emph{bueno / buena}). You'll watch it in the next few lessons: \emph{bon jour} (masc.) but \emph{bonne nuit} (fem.).
\end{grammarlens}

\subsection*{Your turn}
Say these aloud:

\begin{itemize}
\raggedright
  \item \textquotedblleft{}bon\textquotedblright{} (nasal \emph{bõ}) then \textquotedblleft{}bonne\textquotedblright{} (\emph{bun}) — hear the difference
  \item the four shapes — bon, bonne, bons, bonnes
\end{itemize}

\begin{tcolorbox}[breakable,colback=teal!4,colframe=teal!35!black,title={Before you move on}]
Why do \emph{bon} and \emph{bonne} sound different? (The double \emph{nn} in \emph{bonne} cancels the nasal — \emph{bõ} vs \emph{bun}.) Which Spanish adjective shares \emph{bon}'s Latin root? (\emph{bueno}, from the same \emph{bonus}.)
\end{tcolorbox}
\section[le / la / les]{le, la, les — \textquotedblleft{}the,\textquotedblright{} and the gender it carries}

\label{lesson:FR-C01-le-la}



\begin{quote}
Before your first noun: the words for \textquotedblleft{}the.\textquotedblright{} French splits \textquotedblleft{}the\textquotedblright{} by gender — masculine \emph{le}, feminine \emph{la} — a habit it inherited from Latin (its sibling Spanish did the same, from the same source).
\end{quote}

\begin{sounds}
\begin{itemize}
\raggedright
  \item \textbf{le} = \emph{luh}, \textbf{la} = \emph{lah}, \textbf{les} = \emph{lay} (the \emph{s} is silent — \texttt{silent-final}). $\to$ reference
\end{itemize}
\end{sounds}

\begin{cousinweb}
\textbf{le} — \textquotedblleft{}the\textquotedblright{} (masculine). \textbf{la} — \textquotedblleft{}the\textquotedblright{} (feminine). \textbf{les} — \textquotedblleft{}the\textquotedblright{} (plural, both genders). Roots: Latin demonstratives \textbf{ille / illa / illos} (\textquotedblleft{}that one / those\textquotedblright{}):

\begin{itemize}
\raggedright
  \item \textbf{le} $\leftarrow$ \emph{ille}, \textbf{la} $\leftarrow$ \emph{illa}, \textbf{les} $\leftarrow$ \emph{illos} — the pointing force (\textquotedblleft{}that\textquotedblright{}) faded into a plain \textquotedblleft{}the,\textquotedblright{} and the initial \emph{il-} wore off.
\end{itemize}

\begin{quote}
Spanish cousin, exactly: from the \emph{same ille / illa / illos}, Spanish made \emph{el / la / los}. French kept \emph{le} where Spanish kept \emph{el} — same parent. And those same Latin demonstratives also gave French \textbf{il} (\textquotedblleft{}he\textquotedblright{}) and \textbf{elle} (\textquotedblleft{}she\textquotedblright{}) (Spanish \emph{él / ella}): article and pronoun, one Latin root. (English \textquotedblleft{}the\textquotedblright{} is Germanic — unrelated.)
\end{quote}
\end{cousinweb}

\begin{grammarlens}[title={Grammar Lens: French gender, also from Latin}]
Every French noun is \textbf{masculine or feminine} — a grammatical gender baked into the noun, inherited from Latin: its three genders (masculine, feminine, neuter) collapsed to two, the neuters folding into the masculine. (Spanish, the other Latin daughter, did the very same.) You learn each noun \textbf{with its article} (\emph{le} or \emph{la}), because the gender can't be guessed. French has one convenience here: in the plural, both genders share a single article, \emph{les}.
\end{grammarlens}

\subsection*{Your turn}
Say these aloud:

\begin{itemize}
\raggedright
  \item \textquotedblleft{}le\textquotedblright{}, \textquotedblleft{}la\textquotedblright{}, \textquotedblleft{}les\textquotedblright{}
  \item \textquotedblleft{}le jour\textquotedblright{} (masc.), \textquotedblleft{}la nuit\textquotedblright{} (fem.) — the nouns are coming
\end{itemize}

\begin{tcolorbox}[breakable,colback=teal!4,colframe=teal!35!black,title={Before you move on}]
\emph{le} and \emph{la} come from which Latin words? (\emph{ille} / \emph{illa}, \textquotedblleft{}that one\textquotedblright{} — same as Spanish \emph{el/la}.) What's the single plural \textquotedblleft{}the\textquotedblright{}? (\emph{les}.)
\end{tcolorbox}
\section[jour]{jour — \textquotedblleft{}day,\textquotedblright{} and a detour through Latin that explains English \emph{journey}}

\label{lesson:FR-C01-jour}



\begin{quote}
Your first French noun. In Spanish \textquotedblleft{}day\textquotedblright{} is \emph{día}; in French it's \emph{jour} — they look nothing alike, yet they're cousins. The gap between them is the whole lesson, and it cracks open a shelf of English words.
\end{quote}

\subsection*{What to know first}
\begin{itemize}
\raggedright
  \item le / la / les — gender and the articles; \emph{jour} is masculine, \emph{le jour}.
\end{itemize}

\begin{sounds}
\begin{itemize}
\raggedright
  \item \texttt{j-zh} — \textbf{j} is the soft \emph{zh} of \emph{measure}: \emph{jour} starts \emph{zh-}.
  \item \texttt{vowel-ou} — \textbf{ou} = \emph{oo}. So \emph{jour} = \emph{zhoor}. $\to$ reference
\end{itemize}
\end{sounds}

\begin{cousinweb}
\textbf{le jour} — \textquotedblleft{}the day\textquotedblright{} (masculine). Root: Latin \textbf{diurnum} (\textquotedblleft{}daytime, daily\textquotedblright{}) — itself from \textbf{dies} (\textquotedblleft{}day\textquotedblright{}). That extra step through \emph{diurnum} is why \emph{jour} looks so unlike Spanish \emph{día}:

\begin{itemize}
\raggedright
  \item Spanish took \emph{dies} $\to$ \emph{día} (short path).
  \item French took the derivative \emph{diurnum} $\to$ \emph{jorn} $\to$ \emph{jour} (the \emph{di-} crushed down to \emph{j-}).
\end{itemize}

Same ultimate root, two different Latin words, two very different results. And \emph{diurnum} is a goldmine of English cousins — most of them arrived through French:

\begin{itemize}
\raggedright
  \item \textbf{jour}nal (a \emph{daily} record), \textbf{jour}ney (originally a \emph{day's} travel or work — French \emph{journée}, \textquotedblleft{}a day's worth\textquotedblright{}), \textbf{jour}nalist.
  \item so\textbf{journ} (\textquotedblleft{}to stay {[}for days{]}\textquotedblright{}), ad\textbf{journ} (\emph{à jour}, \textquotedblleft{}to {[}another{]} day\textquotedblright{}).
  \item and straight from \emph{dies}: \textbf{di}urnal, \textbf{di}ary — cousins of \emph{jour} through the grandparent root.
\end{itemize}

\begin{quote}
So English \emph{journey} and Spanish \emph{día} are cousins, meeting back at Latin \emph{dies} — one through French \emph{jour}, one direct. That's the kind of hidden family this whole book is built to show.
\end{quote}
\end{cousinweb}

\begin{grammarlens}[title={Grammar Lens}]
\emph{jour} is \textbf{masculine} — \emph{le jour} — a gender it inherited from its Latin source. As always, learn it with the article. Plural: French adds a (silent) \textbf{-s} — \emph{les jours} — sounds identical to the singular; only \emph{les} tells your ear it's plural.
\end{grammarlens}

\subsection*{Your turn}
Say these aloud:

\begin{itemize}
\raggedright
  \item \textquotedblleft{}le jour\textquotedblright{} — \emph{luh zhoor}
  \item \textquotedblleft{}jour\textquotedblright{} then \textquotedblleft{}journal\textquotedblright{}, \textquotedblleft{}journey\textquotedblright{} — the daily root
  \item \textquotedblleft{}les jours\textquotedblright{} — plural, heard only in the \emph{les}
\end{itemize}

\begin{tcolorbox}[breakable,colback=teal!4,colframe=teal!35!black,title={Before you move on}]
\emph{jour} and Spanish \emph{día} share which grandparent root? (Latin \emph{dies}.) What English travel-word is literally \textquotedblleft{}a day's worth\textquotedblright{}? (\emph{Journey}, from \emph{journée}.)
\end{tcolorbox}
\section[bonjour]{bonjour — assembled, and \emph{not} plural (unlike Spanish)}

\label{lesson:FR-C01-bonjour}



\begin{quote}
Snap two words together — \emph{bon} + \emph{jour} — and you have the single most important word in French politeness. And there's a sharp contrast with Spanish hiding in it.
\end{quote}

\subsection*{What to know first}
\begin{itemize}
\raggedright
  \item bon / bonne · jour
\end{itemize}

\begin{cousinweb}
\textbf{bonjour} = \textbf{bon} (\textquotedblleft{}good\textquotedblright{}) + \textbf{jour} (\textquotedblleft{}day\textquotedblright{}) = \textquotedblleft{}good day.\textquotedblright{} One written word, said \emph{bõ- zhoor}. \emph{bon} is masculine because \emph{jour} is masculine — the agreement from the \emph{bon} lesson, quietly at work.
\end{cousinweb}

\begin{culture}
Here's the contrast worth banking. Spanish says \emph{buenos días} — \textbf{plural}, \textquotedblleft{}good days,\textquotedblright{} the fossil of a blessing (\textquotedblleft{}may God give you good days\textquotedblright{}). French says \emph{bonjour} — \textbf{singular}, just \textquotedblleft{}good day.\textquotedblright{} French didn't fossilize the plural; it kept the plain, single wish. Same two building blocks (\emph{good} + \emph{day}), same Latin roots, but Spanish pluralized the blessing and French didn't. Two neighbours, two instincts.

Culturally, \emph{bonjour} carries weight far beyond \textquotedblleft{}hello\textquotedblright{}: in France, opening almost any interaction — a shop, a bus, a desk — \emph{without} a \emph{bonjour} reads as rude. It's the formal, all-purpose daytime greeting (switch to \emph{bonsoir} in the evening), and the polite default where \emph{salut} would be too casual.
\end{culture}

\subsection*{Your turn}
Say these aloud:

\begin{itemize}
\raggedright
  \item build it — \textquotedblleft{}bon\textquotedblright{} … \textquotedblleft{}jour\textquotedblright{} … \textquotedblleft{}bonjour\textquotedblright{}
  \item \textquotedblleft{}salut\textquotedblright{} (casual) then \textquotedblleft{}bonjour\textquotedblright{} (polite) — feel the register gap
  \item bonjour vs Spanish buenos días — one singular, one plural
\end{itemize}

\begin{tcolorbox}[breakable,colback=teal!4,colframe=teal!35!black,title={Before you move on}]
\emph{bonjour} is singular; the Spanish equivalent is plural — what is it, and why plural? (\emph{Buenos días}, fossil of \textquotedblleft{}may God give you good days.\textquotedblright{})
\end{tcolorbox}
\section[soir]{soir — \textquotedblleft{}evening,\textquotedblright{} and the same \textquotedblleft{}late\textquotedblright{} idea Spanish used}

\label{lesson:FR-C01-soir}



\begin{quote}
The evening counterpart of \emph{jour}. And there's a lovely parallel with Spanish waiting: both languages built their \textquotedblleft{}later in the day\textquotedblright{} word out of the Latin word for \emph{late}.
\end{quote}

\subsection*{What to know first}
\begin{itemize}
\raggedright
  \item jour — the day word; \emph{soir} is its evening partner.
\end{itemize}

\begin{sounds}
\begin{itemize}
\raggedright
  \item \texttt{vowel-oi} — \textbf{oi} = \emph{wa}: \emph{soir} is \emph{swahr}. $\to$ reference
\end{itemize}
\end{sounds}

\begin{cousinweb}
\textbf{le soir} — \textquotedblleft{}the evening\textquotedblright{} (masculine). Root: Latin \textbf{sērum} (\emph{tempus}), \textquotedblleft{}the \textbf{late} {[}hour{]},\textquotedblright{} from \textbf{sērus} (\textquotedblleft{}late\textquotedblright{}). English cousins are few but real: \emph{soirée} (an evening party — borrowed straight from French \emph{soir}), and the rare \emph{serotine} (a bat that flies in the \emph{late} evening).

\begin{quote}
Spanish parallel — this is the good part. Spanish \emph{tarde} (\textquotedblleft{}afternoon\textquotedblright{}) comes from Latin \emph{tardus}, \textquotedblleft{}\textbf{late}.\textquotedblright{} French \emph{soir} (\textquotedblleft{}evening\textquotedblright{}) comes from \emph{sērus}, \textquotedblleft{}\textbf{late}.\textquotedblright{} Two different Latin \textquotedblleft{}late\textquotedblright{} words, two languages, the same idea: name the later part of the day after its \emph{lateness}.
\end{quote}
\end{cousinweb}

\begin{grammarlens}[title={Grammar Lens}]
\emph{soir} is \textbf{masculine} — \emph{le soir} — so, like \emph{jour}, it will take \emph{bon} (not \emph{bonne}) in the greeting. Plural \emph{les soirs}, the \emph{s} silent.
\end{grammarlens}

\subsection*{Your turn}
Say these aloud:

\begin{itemize}
\raggedright
  \item \textquotedblleft{}le soir\textquotedblright{} — \emph{luh swahr}
  \item \textquotedblleft{}soir\textquotedblright{} then English \textquotedblleft{}soirée\textquotedblright{}
  \item Spanish \emph{tarde} $\leftarrow$ \textquotedblleft{}late\textquotedblright{}, French \emph{soir} $\leftarrow$ \textquotedblleft{}late\textquotedblright{} — same idea
\end{itemize}

\begin{tcolorbox}[breakable,colback=teal!4,colframe=teal!35!black,title={Before you move on}]
What did \emph{soir}'s Latin root \emph{sērus} mean? (\textquotedblleft{}Late.\textquotedblright{}) What Spanish time-word did the exact same thing from a different \textquotedblleft{}late\textquotedblright{} word? (\emph{Tarde}, $\leftarrow$ \emph{tardus}.)
\end{tcolorbox}
\section[bonsoir]{bonsoir — good evening}

\label{lesson:FR-C01-bonsoir}



\begin{quote}
Same assembly as \emph{bonjour}, one word swapped. If you have that one, this is nearly free.
\end{quote}

\subsection*{What to know first}
\begin{itemize}
\raggedright
  \item bon / bonne · soir · bonjour
\end{itemize}

\begin{cousinweb}
\textbf{bonsoir} = \textbf{bon} (\textquotedblleft{}good\textquotedblright{}) + \textbf{soir} (\textquotedblleft{}evening\textquotedblright{}) = \textquotedblleft{}good evening.\textquotedblright{} \emph{bon} stays masculine because \emph{soir} is masculine — same as \emph{bonjour}, no new agreement to learn. Said \emph{bõ- swahr}.
\end{cousinweb}

\begin{culture}
\emph{bonsoir} is the evening swap for \emph{bonjour}: as the day turns to evening (roughly late afternoon on), French speakers switch from \emph{bonjour} to \emph{bonsoir}. Both are the polite, all-purpose greeting for their time of day. And like Spanish \emph{buenas noches}, \emph{bonsoir} works both \textbf{arriving} and \textbf{leaving} in the evening — \textquotedblleft{}good evening\textquotedblright{} on the way in, \textquotedblleft{}goodnight-ish\textquotedblright{} on the way out. (For actually sending someone to bed, French has a separate phrase — the next two lessons.)
\end{culture}

\subsection*{Your turn}
Say these aloud:

\begin{itemize}
\raggedright
  \item \textquotedblleft{}bonsoir\textquotedblright{} — \emph{bõ- swahr}
  \item \textquotedblleft{}bonjour\textquotedblright{} (daytime) then \textquotedblleft{}bonsoir\textquotedblright{} (evening) — the pair
\end{itemize}

\begin{tcolorbox}[breakable,colback=teal!4,colframe=teal!35!black,title={Before you move on}]
What two words build \emph{bonsoir}, and why does \emph{bon} stay masculine? (\emph{bon} + \emph{soir}; \emph{soir} is masculine.)
\end{tcolorbox}
\section[nuit]{nuit — \textquotedblleft{}night,\textquotedblright{} and a sound-change that mirrors Spanish}

\label{lesson:FR-C01-nuit}



\begin{quote}
\textquotedblleft{}Night.\textquotedblright{} In Spanish it's \emph{noche}; in French, \emph{nuit}. They came from the \emph{same} Latin word and drifted apart by two different rules — and once you see the rules, you can predict a dozen more words.
\end{quote}

\subsection*{What to know first}
\begin{itemize}
\raggedright
  \item le / la / les — \emph{nuit} is feminine, \emph{la nuit}.
\end{itemize}

\begin{sounds}
\begin{itemize}
\raggedright
  \item \textbf{nuit} = \emph{nwee} — the \textbf{u} is the front-rounded French \emph{u} gliding into \emph{ee}; the final \textbf{t} is silent (\texttt{silent-final}). $\to$ reference
\end{itemize}
\end{sounds}

\begin{cousinweb}
\textbf{la nuit} — \textquotedblleft{}the night\textquotedblright{} (feminine). Root: Latin \textbf{noctem} (\emph{nox}, stem \emph{noct-}). English cousins through the \emph{noct-} root: \textbf{noct}urnal, \textbf{noct}urne, equi\textbf{nox} (\textquotedblleft{}equal night\textquotedblright{}).

\textbf{The sound-change, side by side with Spanish.} Latin \emph{noctem} went two ways:

\begin{itemize}
\raggedright
  \item Spanish softened \emph{-ct-} $\to$ \textbf{-ch-}: \emph{noctem $\to$ noche}.
  \item French turned \emph{-ct-} $\to$ \textbf{-it-}: \emph{noctem $\to$ nuit}.
\end{itemize}

And it's a \emph{rule}, not a one-off — watch it repeat:

\noindent
\begin{tabularx}{\linewidth}{@{}>{\raggedright\arraybackslash}X>{\raggedright\arraybackslash}X>{\raggedright\arraybackslash}X>{\raggedright\arraybackslash}X@{}}
\toprule
\textbf{Latin} & \textbf{Spanish} & \textbf{French} & \textbf{English} \\
\midrule
\emph{noctem} & noche & nuit & night \\
\emph{lactem} & leche & lait & (lact-, \textquotedblleft{}milk\textquotedblright{}) \\
\emph{factum} & hecho & fait & fact \\
\bottomrule
\end{tabularx}

So a Spanish \emph{-ch-} word and a French \emph{-it-} word are often the \emph{same} Latin word — and the English cousin usually kept the \emph{-ct-}.
\end{cousinweb}

\begin{grammarlens}[title={Grammar Lens}]
\emph{nuit} is \textbf{feminine} — \emph{la nuit} — inherited from Latin, where \emph{nox} was feminine (exactly as Spanish \emph{la noche} is). Its gender matters next lesson: it's why \textquotedblleft{}good night\textquotedblright{} is \emph{bon\textbf{ne} nuit}, not \emph{bon nuit}.
\end{grammarlens}

\subsection*{Your turn}
Say these aloud:

\begin{itemize}
\raggedright
  \item \textquotedblleft{}la nuit\textquotedblright{} — \emph{lah nwee}
  \item \textquotedblleft{}nuit\textquotedblright{} then Spanish \textquotedblleft{}noche\textquotedblright{} — same Latin \emph{noctem}, two rules
  \item the row — \textquotedblleft{}noche, nuit, night\textquotedblright{}
\end{itemize}

\begin{tcolorbox}[breakable,colback=teal!4,colframe=teal!35!black,title={Before you move on}]
Latin \emph{noctem} $\to$ Spanish \emph{noche} by \emph{-ct-$\to$-ch-}; $\to$ French \emph{nuit} by which change? (\emph{-ct-$\to$-it-}.) Is \emph{nuit} masculine or feminine? (Feminine — \emph{la nuit}.)
\end{tcolorbox}
\section[bonne nuit]{bonne nuit — good night, and the feminine side of \emph{bon}}

\label{lesson:FR-C01-bonne-nuit}



\begin{quote}
The last greeting of the day — and the one that finally shows \emph{bon} in its feminine coat.
\end{quote}

\subsection*{What to know first}
\begin{itemize}
\raggedright
  \item bon / bonne · nuit
\end{itemize}

\begin{cousinweb}
\textbf{bonne nuit} = \textbf{bonne} (\textquotedblleft{}good,\textquotedblright{} feminine) + \textbf{nuit} (\textquotedblleft{}night,\textquotedblright{} feminine) = \textquotedblleft{}good night.\textquotedblright{} Here it is at last: because \emph{nuit} is \textbf{feminine}, \emph{bon} becomes \textbf{bonne} — and, from the \emph{bon} lesson, remember they even \emph{sound} different: \emph{bon} is nasal \emph{bõ}, but \emph{bonne nuit} is \emph{bun nwee}, with a real \emph{n}.

So the three evening-ish greetings line up and show the whole agreement rule in one breath:

\begin{itemize}
\raggedright
  \item \emph{bon\textbf{jour}} — \emph{jour} masculine $\to$ \emph{bon}
  \item \emph{bon\textbf{soir}} — \emph{soir} masculine $\to$ \emph{bon}
  \item \emph{bon\textbf{ne} nuit} — \emph{nuit} feminine $\to$ \emph{bonne}
\end{itemize}
\end{cousinweb}

\begin{culture}
Unlike \emph{bonsoir} (an evening \emph{greeting}), \emph{bonne nuit} is a \textbf{farewell for bed} — \textquotedblleft{}goodnight,\textquotedblright{} said when parting for sleep. It maps onto the Spanish split you may feel forming: Spanish overloads \emph{buenas noches} for both, while French keeps \emph{bonsoir} (evening hello/bye) separate from \emph{bonne nuit} (off-to-sleep).
\end{culture}

\subsection*{Your turn}
Say these aloud:

\begin{itemize}
\raggedright
  \item \textquotedblleft{}bonne nuit\textquotedblright{} — \emph{bun nwee}
  \item the trio — bonjour, bonsoir, bonne nuit — and for each, is \emph{bon} masculine or feminine, and why?
\end{itemize}

\begin{tcolorbox}[breakable,colback=teal!4,colframe=teal!35!black,title={Before you move on}]
Why is it \emph{bonne} nuit but \emph{bon} jour? (\emph{nuit} is feminine, \emph{jour} masculine — agreement.) \emph{bonsoir} vs \emph{bonne nuit} — which is \textquotedblleft{}off to bed\textquotedblright{}? (\emph{Bonne nuit}.)
\end{tcolorbox}
\section[Practice]{Practice — the French greetings}

\label{lesson:FR-C01-practice}



\begin{quote}
No new words — just the ones you built, until the register and the agreement are automatic.
\end{quote}

\begin{grammarlens}[title={What you've built}]
\begin{itemize}
\raggedright
  \item \textbf{salut} — casual \textquotedblleft{}hi/bye\textquotedblright{} ($\leftarrow$ \emph{salus}, \textquotedblleft{}health\textquotedblright{}).
  \item \textbf{bien} — \textquotedblleft{}well,\textquotedblright{} a one-word answer ($\leftarrow$ \emph{bene}; Spanish kept the same word).
  \item \textbf{bon / bonne} — \textquotedblleft{}good,\textquotedblright{} the adjective ($\leftarrow$ \emph{bonus}), and its agreement.
  \item \textbf{le / la / les} — \textquotedblleft{}the,\textquotedblright{} carrying gender ($\leftarrow$ \emph{ille/illa/illos}).
  \item \textbf{bonjour} = bon + jour (masc.) — polite daytime \textquotedblleft{}hello,\textquotedblright{} \emph{singular} where Spanish is plural.
  \item \textbf{bonsoir} = bon + soir (masc.) — evening greeting.
  \item \textbf{bonne nuit} = bonne + nuit (fem.) — off-to-bed \textquotedblleft{}goodnight.\textquotedblright{}
\end{itemize}

Two threads run through it all: \textbf{register} (\emph{salut} vs \emph{bonjour}) and \textbf{agreement} (\emph{bon jour} / \emph{bon soir} / \emph{bonne nuit}).
\end{grammarlens}

\subsection*{Your turn}
Say these aloud:

\begin{itemize}
\raggedright
  \item greet a friend casually, then a shopkeeper politely — \emph{salut} / \emph{bonjour}
  \item pick the right greeting for right now (day $\to$ \emph{bonjour}, evening $\to$ \emph{bonsoir})
  \item bonjour, bonsoir, bonne nuit — and say why each takes \emph{bon} or \emph{bonne}
  \item answer \textquotedblleft{}how are you\textquotedblright{} in one word — \emph{bien}, then \emph{très bien}
\end{itemize}

\emph{Twice through:} Run bonjour $\to$ bonsoir $\to$ bonne nuit without hesitating on \emph{bon}/\emph{bonne}.

\subsection*{Writing — close the first runway}
Cover the word \textbf{salut} elsewhere in this chapter. Write the casual greeting once from memory, then uncover it and check the silent final \textbf{t}. This is a chapter-level retrieval of the tiny trace-to-dictation runway, not a new word.

\begin{tcolorbox}[breakable,colback=teal!4,colframe=teal!35!black,title={Before you move on}]
French \emph{bonjour} is singular; the Spanish morning greeting is plural — name it. (\emph{Buenos días}.) Next chapter: introducing yourself — \emph{je m'appelle…} — and French's own formal/informal \textquotedblleft{}you,\textquotedblright{} \textbf{tu} vs \textbf{vous}.
\end{tcolorbox}
